\section{Limitations and conclusion}
\label{sec:conclusion}

\paragraph{Limitations.} The building is a stylised one: uniform floor
population, a single lobby, no express zones or sky lobbies, no time-of-day
peaks, and a resident whose destinations are uniform over the other floors.
The simulator implements collective control with nearest-car dispatch, not a
modern destination-dispatch group controller, and its parking policies are
the three textbook ones; real controllers learn where to park, which would
move our ``lobby'' and ``random'' limits toward each other. The kinematics are
mid-rise values (\KinV\,m/s), so the floor-by-floor run model coincides with
the trapezoid; high-speed cars would need the exact profile, which the code
provides. The stairs rule is deterministic on expectations and ignores effort,
carrying things, and the building's actual stair placement. The resident's
probing trips are frequent enough to perturb a loaded single car slightly.
All of these are configuration knobs, and the accompanying full-scale
configuration widens every grid.

\paragraph{Conclusion.} The middle floor is the right answer to a question
nobody asks --- where to be if every trip is to a random floor and the
elevator is waiting outside your door. Once the elevator has to come and get
you, the answer depends on where it waits; once most trips go to the lobby,
the answer is the lowest floor you can get, by a margin that grows with the
building and with the number of cars. Traffic blurs the curve but does not
move the optimum up until the car is saturated, and there the effect that
lifts it is a queueing one (the low floors are served last on the down sweep)
rather than a geometric one. Stairs make the low floors better still and shift
a lobby-free building's optimum upward by absorbing short trips, but the
cost of climbing is not what does it. And the sense in which ``it does not
matter'' is precise: the indifference band is one floor wide wherever lobby
trips are common, and spans most of the building only where they are rare
and the car is busy. Up or down? Down, unless you never leave.
